\section{RF-LR Gram Marchenko–Pastur spectrum for three-layer case}\label{sec:rf-lr-gram-marchenko-pastur-spectrum-for-three-layer-case}

We now study the spectrum of the RF-LR NTK Gram matrix at initialization in a random-matrix scaling regime. This section concerns the \emph{uncentered} Gram matrix and its spike--bulk limit; we also explain how centering removes the rank-one spike and isolates the spectrum on \( \mathbf{1}^\perp \), which is the relevant subspace after centering targets (or for mean-zero targets).
We work under strong scaling and data assumptions for tractability.

\subsection{General setup and scaling}

\begin{assumption}[Random-matrix scaling]\label{assump:RMT}
As the number of data points \(n\to\infty\):
\begin{itemize}
    \item Hidden widths scale as \( N \sim n^2 \) (infinite-width linearization regime).
    \item Rank scales as \( r \asymp n \) with ratio \( \gamma_{\mathrm{ratio}} \equiv \lim_{n\to\infty} \tfrac{n}{r} \in (0,\infty) \) (Marchenko–Pastur regime).
    \item Input dimension scales linearly with rank: \( d \asymp r \) (so \( d \asymp n \) as well).
\end{itemize}
\end{assumption}

We now introduce the random-matrix theory conditions that are necessary for the asymptotic analysis. Those strong assumptions can be removed in future work using techniques from \cite{benigni2025eigenvaluedistributionneuraltangent}

\begin{assumption}[RMT Gaussian data]\label{assump:RMT-conditions}
Following the framework of~\cite{El_Karoui_2010}, we impose the following conditions  as \( d \to \infty \):
\begin{itemize}
    
    \item[(d)] \textbf{Data generation:} Input vectors satisfy \( X_i = \Sigma_d^{1/2} Y_i \), where \( \Sigma_d^{1/2} \) is the positive semi-definite square root of \( \Sigma_d \).
    
    \item[(e)] \textbf{Entrywise conditions:} The entries of \( Y_i \), a \( d \)-dimensional random vector, are i.i.d. Denoting by \( Y_i^{(k)} \) the \( k \)-th entry of \( Y_i \), we assume \( \mathbb{E}[Y_i^{(k)}] = 0 \), \( \mathrm{var}(Y_i^{(k)}) = 1 \), and \( \mathbb{E}[|Y_i^{(k)}|^{4+\varepsilon}] < \infty \) for some \( \varepsilon > 0 \) (i.e., \( Y_i \) has \( 4+\varepsilon \) absolute moments).
    
    \item[(b)] \textbf{Covariance structure:} Let \( \Sigma_d \) be the \( d \times d \) covariance matrix of the input distribution. Then \( \Sigma_d \) is positive semi-definite and its operator norm \( \|\Sigma_d\|_2 = \sigma_1(\Sigma_d) \) remains bounded in \( d \), that is, there exists \( K > 0 \) such that \( \sigma_1(\Sigma_d) \leq K \) for all \( d \).
    
    \item[(c)] \textbf{Trace normalization:} \( \Sigma_d/d \) has a finite limit, that is, there exists \( \ell \in \mathbb{R} \) such that \( \lim_{d\to\infty} \mathrm{trace}(\Sigma_d)/d = \ell \).
    \end{itemize}
\end{assumption}

With these scaling assumptions in place, we can now analyze the microscopic structure of the kernel. The key insight is that the randomness in the low-rank projections decomposes into independent angular and radial components.

To apply \cite{El_Karoui_2010}, the kernel function \( K_{\infty}(\rho) = \Theta^{(1)}(\rho)\big(1-\arccos(\rho)/\pi\big) + \Sigma^{(1)}(\rho) + 1 \) must be \( C^1 \) in a neighborhood of \( \ell = \lim_{d\to\infty} \mathrm{trace}(\Sigma_d)/d \) and \( C^3 \) in a neighborhood of \( 0 \). This holds trivially since \( \arccos(\rho) \) is \( C^\infty \) on \( (-1,1) \) and the ReLU kernels \( \Theta^{(1)}(\rho) \) and \( \Sigma^{(1)}(\rho) \) have the required smoothness properties.

\subsection{Signal–noise decomposition of the empirical NTK}
\begin{theorem}[Asymptotic normality of the low-rank kernel]\label{thm:kernel_clt}
Let \(x, y \in \mathbb{R}^r\) be rank-\(r\) Gaussian projections with population correlation \( \rho \). Define
\[
\Psi(u,w) \;=\; \Theta^{(1)}(u)\!\left(1-\frac{\arccos(u)}{\pi}\right) \;+\; w\,\Sigma^{(1)}(u) \;+\; 1,
\]
so that \( \hat{\Theta}^{(2)}_r = 1 + \frac{1}{r}\big(\Psi(\hat{\rho}_r, \hat{w}_r) - 1\big) \) from the \( 1/r \) recursion, where \( \hat{\rho}_r \) is the sample correlation and \( \hat{w}_r = \|x\|\,\|y\|/r \). Define the deterministic limit \( K_{\infty}(\rho) = \Psi(\rho,1) \). Under Assumption~\ref{assump:RMT}, as \( r\to\infty \),
\[
\hat{\Theta}^{(2)}_r \;=\; \Theta^{(2)}_{\mathrm{obs}}(\rho) \;\xrightarrow{d}\; K_{\infty}(\rho) \;+\; \frac{1}{\sqrt{r}}\,\mathcal{G}(\rho) + o(r^{-1/2}),
\]
where \( K_{\infty}(\rho)=\mathbb{E}[\Theta^{(2)}_{\mathrm{obs}}(\rho)] \) is the deterministic limit and \( \mathcal{G}(\rho) \sim \mathcal{N}\!\big(0,\sigma_{\mathrm{tot}}^2(\rho)\big) \) is a centered Gaussian field with variance
\[
\sigma_{\mathrm{tot}}^2(\rho) \;=\; \mathrm{Var}(\mathcal{G}(\rho)) \;=\; \underbrace{\big[\partial_\rho K_{\infty}(\rho)\,(1-\rho^2)\big]^2}_{\text{Fisher (angle)~\cite{fisher1915probable,hotelling1953new}}} \;+\; \underbrace{\big[\Sigma^{(1)}(\rho)\,\sqrt{1+\rho^2}\big]^2}_{\text{Kibble (norms)}}.
\]
\end{theorem}

\paragraph{Proof intuition.}
By Fisher–Kibble decoupling (Lemma~\ref{lem:fisher-kibble})~\cite{fisher1915probable,hotelling1953new}, the angular \( \hat{\rho}_r \) and radial \( \hat{w}_r \) components are independent. A joint CLT gives \( \sqrt{r}\,(\hat{\rho}_r-\rho,\ \hat{w}_r-1) \xrightarrow{d} \mathcal{N}(\mathbf{0},\ \mathrm{diag}((1-\rho^2)^2,\ 1+\rho^2)) \). Applying the Delta method to \( \Psi \) at \( (\rho,1) \) with gradient \( (\partial_\rho K_{\infty}(\rho),\ \Sigma^{(1)}(\rho)) \) yields the stated variance (no cross term by independence). The limit is a deterministic kernel plus \( O_p(r^{-1/2}) \) Gaussian fluctuation.

Proof: See Appendix~\ref{appendix:proof-kernel-clt}.

Having established the entrywise asymptotic distribution, we now turn to the macroscopic spectral properties of the kernel matrix. The eigenvalue distribution follows a deformed Marchenko–Pastur law with a characteristic spike-bulk structure.

\subsection{Macroscopic spectrum: deformed Marchenko–Pastur law}
Recent work by~\cite{benigni2025eigenvaluedistributionneuraltangent} establishes the eigenvalue distribution of the NTK in the quadratic scaling regime (\( N \sim n^2 \)) for single-layer networks, showing convergence to a deformed Marchenko–Pastur law with a spike-bulk structure. We extend this analysis to the three-layer case with random features and rank-\( r \) bottlenecks, where the frozen feature weights and low-rank structure modify both the spike location and bulk support.

\paragraph{Centered vs uncentered Gram matrices (removing the rank-one spike).}
Let \(K\in\mathbb{R}^{n\times n}\) be a symmetric Gram matrix built from a kernel evaluated on a dataset \(\{x_i\}_{i=1}^n\). Define the centering matrix \(H = I_n - \frac{1}{n}\mathbf{1}\mathbf{1}^\top\) and the centered Gram matrix \(K_c = HKH\). Then for any \(v\perp \mathbf{1}\) we have \(Hv=v\) and therefore
\[
v^\top K_c v = v^\top HKH v = v^\top K v.
\]
Consequently, the eigenvalues of \(K\) restricted to the subspace \(\mathbf{1}^\perp\) coincide with the (nonzero) eigenvalues of \(K_c\). In particular, if \(K\) contains a rank-one component aligned with \(\mathbf{1}\) (e.g. arising from an additive constant term in the kernel), this produces a large top eigenvalue in the uncentered spectrum but is removed by centering. For kernel regression with mean-zero targets (or after centering labels), convergence rates are governed by the spectrum on \(\mathbf{1}^\perp\).

\begin{theorem}[Spike + bulk spectral limit]\label{thm:spike-bulk}
Let \( \mathbf{M} = [\Theta^{(2)}_{ij}]_{i,j=1}^n \). Under Assumptions~\ref{assump:RMT} and~\ref{assump:RMT-conditions}, as \( n\to\infty \), the empirical spectrum of \( \mathbf{M} \) converges to a deformed Marchenko–Pastur law with:
\begin{itemize}
    \item A rank-one outlier (spike) of order \(O(n)\): \( \lambda_{\mathrm{spike}} \approx n\,\alpha \), with \( \alpha = K_{\infty}(0) \).
    \item A continuous bulk of order \(O(1)\) supported on
    \[
    \Big[\;\gamma + \beta\,(1-\sqrt{\gamma_{\mathrm{ratio}}})^2,\;\; \gamma + \beta\,(1+\sqrt{\gamma_{\mathrm{ratio}}})^2\;\Big],
    \]
    where \( \beta = K_{\infty}'(0) \) and \( \gamma = K_{\infty}(\ell) - \alpha - \ell\,\beta \) (diagonal shift; \(\ell\) denotes the kernel's diagonal argument).
\end{itemize}
Proof: See Appendix~\ref{appendix:proof-spike-bulk} (via~\cite{El_Karoui_2010}). This theorem is stated for the deterministic inner-product kernel matrix built from the limiting kernel function \(K_\infty\); incorporating finite-\(r\) Fisher--Kibble fluctuations or finite-width effects into the limiting ESD would require BP-style resolvent/Gaussian-equivalence machinery and is left for future work.
\end{theorem}

\begin{remark}[How to interpret the spike for learning]
The spike in Theorem~\ref{thm:spike-bulk} is a statement about the \emph{uncentered} Gram matrix. When the kernel includes an additive constant term (here the \(+1\) in \(K_\infty\)), the associated rank-one component aligns with \(\mathbf{1}\) and yields \(\lambda_{\max}(\mathbf{M})=\Theta(n)\). This does not by itself imply poor conditioning of the learning problem: after centering (or for mean-zero labels), the relevant spectrum is that of \(H\mathbf{M}H\) on \(\mathbf{1}^\perp\), i.e. the bulk/shifted spectrum.
\end{remark}

\begin{remark}[Bulk laws vs extreme eigenvalues]\label{rmk:bulk-vs-extremes}
Theorem~\ref{thm:spike-bulk} is a statement about the \emph{empirical spectral distribution} (a bulk law) under strong random-matrix assumptions. Even when a limiting bulk measure exists and has compact support, this does not automatically control the \emph{extreme} eigenvalues at finite \(n\): the smallest eigenvalue can be sensitive to (i) finite-size edge fluctuations; (ii) deviations from the independence/isotropy assumptions; and (iii) additional structured components beyond the rank-one spike (e.g. heteroskedastic norms, anisotropic covariance \(\Sigma_d\), or higher-order kernel terms).

Centering removes only the rank-one component aligned with \(\mathbf{1}\) and leaves the spectrum on \(\mathbf{1}^\perp\) unchanged (paragraph ``Centered vs uncentered''); it does not make the remaining eigenvalue distribution ``less skewed''. In particular, in depth-scaling regimes where the relevant centered scale becomes small (e.g. \(\lambda_{\min}(HK^{(L)}H)\asymp 1/(rL)\) in the equicorrelated proxy of Section~\ref{sec:depth-scaling}), random fluctuations can drive \(\lambda_{\min}(H\widehat{K}H)\) very close to \(0\) even though \(\widehat{K}\succeq 0\) always holds (Remark~\ref{rmk:lambda-min-negative-appendix} in Appendix~\ref{appendix:lambda-min-negative}).

Controlling the smallest eigenvalue beyond bulk limits typically requires separate ``local law'' or edge universality inputs, which are outside the scope of the present random-matrix section.
\end{remark}

We can summarize the three-layer results as follows.

\paragraph{three-layer case summary.}
Specializing to two layers with ReLU features, the deterministic kernel limit is
\( K_{\infty}(\rho) = \Theta^{(1)}(\rho)\big(1-\arccos(\rho)/\pi\big) + \Sigma^{(1)}(\rho) + 1 \).
At the dataset level, the corresponding Gram matrix exhibits a rank-one ``spike'' of size
\( \lambda_{\mathrm{spike}} \approx n\,K_{\infty}(0) \), superimposed on a deformed Marchenko–Pastur bulk.
The bulk support is the interval given by Theorem~\ref{thm:spike-bulk}, with slope
\( \beta = K'_{\infty}(0) \) and diagonal shift \( \gamma \) determined by the kernel's on-diagonal value.
